% =============================================================================
% Kapitel 5: Prototypische Realisierung
% =============================================================================

\section{Prototypische Realisierung}
\label{sec:prototype_implementation}


% =============================================================================
\subsection{Realisierungsumfang und technische Basis}
\label{sec:prototype_scope_stack}
% =============================================================================

Kapitel~4 beschreibt die konzipierte Informationsverarbeitungsarchitektur des
Turing Generators unabhängig von ihrer konkreten technischen Umsetzung.
Im Folgenden wird dargestellt, wie die für die Forschungsfrage wesentlichen
Architekturverantwortlichkeiten in einem lauffähigen Proof of Concept
realisiert wurden.

Die Darstellung folgt dabei nicht der internen Paket- oder Modulstruktur der
Software, sondern dem implementierten Engineering Information Flow. Ausgangspunkt
ist eine registrierte Engineering Source. Anschließend wird betrachtet, welche
technischen Artefakte und Verarbeitungsschritte entstehen, wie probabilistische
Interpretation und menschliche Entscheidungen voneinander getrennt werden und
wie daraus schrittweise ein validierbares SysML-v2-Artefakt entsteht.

Beschrieben wird der zum Abschluss der prototypischen Entwicklung akzeptierte
Realisierungsstand. Die zeitliche Reihenfolge einzelner Implementierungsschritte,
Architecture Decision Records oder Work Packages bildet daher nicht die
Gliederungslogik dieses Kapitels. Frühere Realisierungsvarianten werden nur dort
aufgegriffen, wo ihre Abgrenzung zum finalen Zustand für das Verständnis einer
wesentlichen technischen Entscheidung erforderlich ist.

Die konkreten Ergebnisse der Verifikation und Validierung werden dabei nicht
vorweggenommen. Kapitel~5 beschreibt zunächst den implementierten Zustand,
während Kapitel~6 die zugehörigen Untersuchungen und Kapitel~7 deren Ergebnisse
behandelt.


\subsubsection{Abgrenzung zwischen Target Architecture und Proof of Concept}
\label{sec:prototype_scope}

Wie bereits in Abschnitt~\ref{sec:target_vs_mvp} abgegrenzt, ist die in
Kapitel~4 beschriebene Target Architecture nicht mit dem vollständig
implementierten Funktionsumfang des Proof of Concept gleichzusetzen.
Das Architekturmodell beschreibt die vorgesehenen Verantwortlichkeiten,
Informationsgrenzen und Authority Boundaries des Turing Generators, während
der Proof of Concept einen für die Untersuchung ausgewählten Ausschnitt davon
technisch realisiert.

Im Mittelpunkt der prototypischen Umsetzung steht die thesis-relevante
End-to-End-Verarbeitungskette. Diese reicht von der Registrierung und
Aufbereitung einer Engineering Source über die source-grounded
Informationsverarbeitung, Human Engineering Review und projektweite
Modellableitung bis zur Generierung, Validierung und finalen menschlichen
Freigabe einer SysML-v2-Repräsentation.

Damit werden insbesondere diejenigen Architekturmechanismen umgesetzt, die
für die Untersuchung der kontrollierten Überführung heterogener
Engineering-Information entscheidend sind. Dazu gehören die explizite
Bindung an Project- und Source-Kontexte, die Trennung von Evidence und
Interpretation, die Beschränkung maschineller Verarbeitung auf nicht
autoritative Verarbeitungsschritte, menschliche Authority Boundaries sowie
die nachvollziehbare Ableitung und Generierung des Zielmodells.

Nicht Bestandteil des Proof of Concept ist dagegen eine vollständige
Produktimplementierung sämtlicher in der Target Architecture denkbarer
Betriebs-, Deployment- und Skalierungsaspekte. Insbesondere stellt der
Prototyp keinen Nachweis allgemeiner Production Readiness dar.

Für die weitere Darstellung werden deshalb drei Ebenen auseinandergehalten:
die modellierte Target Architecture, der im Repository realisierte
Proof-of-Concept-Zustand und das innerhalb der Verifikation und Validierung
tatsächlich untersuchte Verhalten. Kapitel~5 bezieht sich ausschließlich auf
die zweite dieser Ebenen.


\subsubsection{Technologiestack und Ausführungsumgebung}
\label{sec:technology_stack}

Der Proof of Concept wurde als Python-basierte Anwendung realisiert.
Die Implementierungsstruktur ist modular aufgebaut und trennt die
Verantwortlichkeiten der einzelnen Verarbeitungsschritte über explizite
Contracts und persistierte Zwischenartefakte voneinander.

Für die Interaktion mit dem System wird eine Streamlit-basierte Oberfläche
verwendet. Sie bildet den Human-facing Zugang zum Engineering Workflow und
stellt insbesondere Projektzustände, Verarbeitungsergebnisse, Evidence sowie
notwendige menschliche Entscheidungen dar. Die Benutzeroberfläche bildet dabei
keine eigene Engineering Authority, sondern stellt die jeweils zulässigen
Interaktionen mit den zugrunde liegenden Verarbeitungs- und
Entscheidungsmechanismen bereit.

Probabilistische Verarbeitungsschritte werden über angebundene
Large-Language-Model-Dienste ausgeführt. LLM-Aufrufe sind jedoch nicht über
den gesamten Verarbeitungsweg verteilt, sondern auf diejenigen Stufen
beschränkt, in denen semantische Interpretation oder die Bewertung
mehrdeutiger Engineering-Information erforderlich ist. Deterministische
Verarbeitung wird dagegen für kontrollierbare Transformationen,
Identitätsbildung, Persistenz, Modellassemblierung, Generierung und
technische Prüfungen verwendet, soweit diese ohne semantische Interpretation
durchgeführt werden können.

Die Verarbeitungsergebnisse werden nicht ausschließlich im Laufzeitzustand
der Anwendung gehalten. Zentrale Informationsobjekte werden als strukturierte
Projektartefakte und Manifeste persistiert. Dadurch können nachfolgende
Verarbeitungsstufen auf explizit identifizierte Eingaben zurückgreifen, statt
impliziten oder lediglich im Arbeitsspeicher vorhandenen Zustand zu verwenden.

\subsubsection{Projektorientierte Persistenz und Artifact Organization}
\label{sec:artifact_organization}


Eine zentrale technische Grundlage des Proof of Concept ist die
projektorientierte Organisation des Verarbeitungszustands. Engineering
Sources, Processing Runs, einzelne Processing Attempts, Human Decisions und
nachgelagerte Modell- und Generierungsartefakte werden nicht als voneinander
unabhängige Dateien behandelt, sondern einem expliziten Projekt- und
Verarbeitungskontext zugeordnet.

Diese Struktur wurde im Entwicklungsverlauf unter anderem durch die
Project-Workspace-Architektur und die Processing-State- und
Artifact-Organization-Entscheidung konkretisiert, die in ADR-005 und
ADR-012 dokumentiert sind. Dadurch wird der technische Zustand einer
Verarbeitung nicht ausschließlich durch den jeweils zuletzt erzeugten Output
repräsentiert. Auch die Herkunft, der Verarbeitungskontext und die
Abhängigkeiten zwischen den Artefakten bleiben erhalten.

Persistierte Manifeste dienen dazu, solche Zustände und Beziehungen explizit
abzubilden. Für relevante Engineering-Artefakte werden stabile technische
Identitäten verwendet. Fingerprints ermöglichen zusätzlich, einen
nachgelagerten Verarbeitungsschritt an den tatsächlich konsumierten
Artefaktzustand zu binden. Ein späterer Human Review oder eine
Modellableitung kann damit nicht lediglich auf einen logischen Namen
verweisen, sondern auf einen eindeutig bestimmbaren technischen
Verarbeitungszustand.

Die Persistenz folgt dabei überwiegend einem immutable beziehungsweise
append-orientierten Prinzip. Bereits dokumentierte Verarbeitungsschritte oder
Entscheidungen werden nicht stillschweigend durch einen späteren Zustand
überschrieben. Stattdessen entstehen neue Artefakte oder Zustände, deren
Beziehung zu vorausgehenden Verarbeitungsschritten nachvollziehbar bleibt.

Diese technische Organisation bildet die Grundlage für die in
Kapitel~4 beschriebene Provenance- und Traceability-Architektur. Sie allein
begründet jedoch noch keine Engineering Authority. Ein persistiertes
Artefakt, eine stabile ID oder ein Fingerprint weist zunächst lediglich einen
bestimmten technischen Zustand nach. Ob dessen Engineering-Inhalt
autorisiert ist, wird erst an den dafür vorgesehenen Human-Authority-Grenzen
entschieden.

Auf dieser Grundlage beginnt der eigentliche implementierte
Informationsverarbeitungsweg mit der Anlage eines Projektkontexts und der
Registrierung einer Engineering Source. Diese erste konkrete Stufe wird im
folgenden Abschnitt betrachtet.

\subsection{Projektanlage, Source-Registrierung und Processing Context}
\label{sec:implementation_context}

Bevor Engineering-Information verarbeitet werden kann, wird sie im
Proof of Concept einem expliziten Projekt- und Source-Kontext zugeordnet.
Damit beginnt die technische Verarbeitung nicht unmittelbar mit dem Inhalt
einer Datei, sondern mit der Herstellung eines eindeutig referenzierbaren
Verarbeitungskontexts.

Diese Bindung ist erforderlich, weil nachgelagerte Artefakte nicht nur auf
ihren jeweiligen Inhalt, sondern auch auf ihre Herkunft und den konkreten
Verarbeitungsvorgang zurückgeführt werden müssen. Project, Source,
Processing Run und Processing Attempt bilden deshalb gemeinsam den
technischen Kontext, innerhalb dessen die nachfolgenden
Informationsverarbeitungsschritte ausgeführt werden.

Die hierfür verwendeten Workspace-, Source-Registry- und
Processing-State-Strukturen wurden unter anderem in ADR-005, ADR-010 und
ADR-012 als Realisierungsentscheidungen dokumentiert. Die zugehörigen
Architecture Decision Records werden entsprechend der in Kapitel~4
eingeführten Systematik im ADR-Anhang zusammengeführt.


\subsubsection{Projektregistrierung und Source-Bindung}
\label{sec:project_source_registration}

Die oberste technische Zuordnungsgrenze bildet das \emph{Project}.
Ein Projekt kapselt den Verarbeitungskontext für die Engineering Sources,
die im Rahmen eines gemeinsamen Modellierungsziels verarbeitet werden sollen.
Dadurch werden Quellen und die daraus entstehenden Artefakte nicht in einem
globalen Verarbeitungskontext geführt, sondern einem expliziten Projekt
zugeordnet.

Engineering Sources werden vor ihrer Verarbeitung registriert und erhalten
eine stabile technische Identität. Neben der eigentlichen Source Information
wird damit ein Source-Kontext erzeugt, der nachgelagerten
Verarbeitungsschritten eine eindeutige Referenz auf die Ausgangsquelle
ermöglicht.

Die Registrierung einer Source stellt dabei noch keine fachliche Bewertung
ihres Inhalts dar. Sie bedeutet insbesondere nicht, dass die enthaltene
Engineering-Information bereits als korrekt, relevant oder für eine spätere
Modellableitung autorisiert gilt. Die Source-Registrierung schafft zunächst
ausschließlich den technischen Kontext, in dem diese Fragen später
bearbeitet werden können.

Zwischen Project und Source besteht deshalb eine explizite Bindung.
Source-bezogene Verarbeitung wird nur innerhalb des ihr zugeordneten
Projektkontexts ausgeführt. Dadurch wird verhindert, dass Artefakte
unterschiedlicher Projekte allein aufgrund gleichartiger Dateinamen,
lokaler IDs oder ähnlicher Inhalte unbeabsichtigt in denselben
Verarbeitungspfad gelangen.

Ein Projekt kann gleichzeitig mehrere registrierte Engineering Sources
enthalten. Diese gemeinsame Projektzugehörigkeit führt jedoch noch nicht zu
einer inhaltlichen Zusammenführung. Die Verarbeitung bleibt zunächst
source-lokal. Erst an der später beschriebenen Project-Fit-Grenze wird
entschieden, welche bereits verarbeitete und autorisierte Information in
die projektweite Modellableitung eingehen darf.

Damit erfüllt die Registrierung zwei unterschiedliche Funktionen:
Das Project bildet die übergeordnete technische Isolationsgrenze,
während die Source Identity die konkrete Herkunft einer
Engineering-Information innerhalb dieses Projekts erhält.


\subsubsection{Processing Run, Attempt und Artifact Identity}
\label{sec:processing_identity}

Die registrierte Source bildet die Grundlage für einen
\emph{Processing Run}. Ein Processing Run beschreibt einen eindeutig
identifizierbaren Verarbeitungskontext für eine konkrete Source innerhalb
eines Projekts. Nachgelagerte Processing Artifacts werden dadurch nicht
lediglich der Ausgangsdatei, sondern auch dem Verarbeitungsvorgang
zugeordnet, in dem sie entstanden sind.

Von einem Processing Run werden einzelne
\emph{Processing Attempts} unterschieden. Dadurch kann derselbe
Verarbeitungsauftrag erneut ausgeführt werden, ohne einen vorherigen
Versuch oder dessen bereits persistierte Artefakte zu überschreiben.
Die unterschiedlichen Attempts bleiben als getrennte technische Zustände
erhalten.

Diese Trennung ist insbesondere für probabilistische Verarbeitungsschritte
relevant. Eine erneute Ausführung muss nicht zwangsläufig denselben
Interpretationsoutput erzeugen. Die Attempt-Bindung ermöglicht daher,
ein konkretes Verarbeitungsergebnis auf genau den Ausführungsversuch
zurückzuführen, aus dem es hervorgegangen ist.

Auch die Identität nachgelagerter Processing Artifacts wird explizit
verwaltet. Artifact IDs werden nicht aus Dateinamen, Anzeigenamen oder
fachlichen Bezeichnungen abgeleitet. Sie dienen als technische Identitäten
innerhalb des jeweiligen Verarbeitungskontexts.

Innerhalb eines Processing Run kann ein Artefakt über die Kombination aus
Artefakttyp und Artifact ID identifiziert werden. Sobald Artefakte mehrerer
Processing Runs in einem gemeinsamen Projektkontext betrachtet werden,
wird zusätzlich die Processing-Run-Identität berücksichtigt. Die
projektweit wirksame technische Referenz folgt damit dem Schema

\[
(\texttt{processing\_run\_id},
 \texttt{artifact\_type},
 \texttt{artifact\_id}).
\]

Gleichlautende lokale Artifact IDs aus unterschiedlichen Processing Runs
bezeichnen dadurch keine identischen Artefakte. Ihre jeweilige
Run-Zugehörigkeit bleibt Bestandteil der projektweiten Identität.

Neben den Identitäten werden für relevante Artefakte Fingerprints und
Manifestinformationen persistiert. Ein nachgelagerter Verarbeitungsschritt
kann dadurch nicht nur auf eine logische Artifact ID, sondern auf den
konkret gebundenen Artefaktzustand verweisen.

Diese technische Bindung besitzt zunächst keine fachliche
Autoritätswirkung. Processing Run, Attempt, Artifact Identity und
Fingerprint dokumentieren, welcher Zustand verarbeitet oder betrachtet
wurde. Ob dessen Engineering-Inhalt später autorisiert werden darf,
entscheidet sich erst an den expliziten Human-Review-Grenzen.

Mit dem erzeugten Project-, Source-, Run- und Attempt-Kontext ist die
Engineering Source eindeutig in den Verarbeitungspfad eingebunden.
Der nächste Schritt besteht darin, ihren Inhalt für die weitere Analyse
aufzubereiten und daraus source-grounded Evidence zu bilden.

\subsection{Source Preparation und Bildung source-grounded Evidence}
\label{sec:implementation_evidence}

Nach der Bindung einer Engineering Source an Project, Processing Run und
Processing Attempt beginnt die inhaltliche Verarbeitung der Source.
Der erste Schritt besteht dabei nicht unmittelbar in einer semantischen
Interpretation durch ein Large Language Model. Stattdessen wird die
Ausgangsinformation zunächst in einen stabilen, source-gebundenen
Verarbeitungszustand überführt.

Die Implementierung trennt hierzu drei aufeinanderfolgende Aufgaben:
die deterministische Aufbereitung der Source Information, die Erkennung
source-grounded Evidence und die Bildung von Engineering Subjects als
fachliche Bezugspunkte für die nachgelagerte Interpretation.

Diese Trennung folgt der in Kapitel~4 beschriebenen Architekturgrenze zwischen
Source Preparation, Evidence und Interpretation. Der zentrale Zweck besteht
darin, die spätere probabilistische Verarbeitung auf explizit identifizierte
und nachvollziehbare Ausgangsinformation zu beziehen.


\subsubsection{Deterministische Source Preparation}
\label{sec:source_preparation}

Die registrierte Engineering Source wird zunächst in eine kontrollierte
interne Repräsentation überführt. Ziel dieses Schritts ist nicht die
fachliche Interpretation ihres Inhalts, sondern die Herstellung eines
stabilen Analysegegenstands für die nachfolgenden Verarbeitungsschritte.

Die Source Preparation ist deshalb deterministisch von der späteren
LLM-basierten Interpretation getrennt. Die Ausgangsinformation wird zunächst in
analysierbare Source Projections überführt, wobei die Beziehung zur
ursprünglichen Source erhalten bleibt.

Innerhalb dieser vorbereiteten Source Information werden anschließend
\emph{Source Analysis Units} gebildet. Eine Source Analysis Unit beschreibt
einen immutable, persona-unabhängigen Analysekontext, auf den sich die
nachfolgenden Evidence- und Interpretationsschritte beziehen können. Sie
grenzt damit den konkreten Source-Bereich ein, ohne bereits eine fachliche
Interpretation seines Inhalts vorzunehmen.

Die resultierenden Verarbeitungseinheiten repräsentieren damit keinen neuen
Engineering-Inhalt. Sie strukturieren die Source Information lediglich so,
dass nachfolgende Verarbeitungsschritte konkrete Bereiche referenzieren
können, ohne die Verbindung zum ursprünglichen Dokument zu verlieren.

Diese Trennung wurde im Projekt unter anderem in ADR-009 als eigenständige
Processing Boundary dokumentiert. Der dort festgelegte Realisierungsansatz
verhindert, dass semantische Interpretation bereits während der technischen
Source-Aufbereitung implizit mitgeführt wird.

Für nachgelagerte Verarbeitungsschritte ist damit eindeutig unterscheidbar,
ob eine Information unmittelbar aus der Source stammt oder erst durch eine
spätere semantische Interpretation entstanden ist.

Diese Unterscheidung bildet die Grundlage für die anschließende
Evidence Detection.


\subsubsection{Source-grounded Evidence Detection}
\label{sec:evidence_detection_impl}

Auf Grundlage der vorbereiteten Source Information wird anschließend
\emph{source-grounded Evidence} gebildet.

Evidence bezeichnet dabei einen eindeutig auf die Engineering Source
zurückführbaren Informationsausschnitt, der für eine spätere fachliche
Interpretation relevant sein kann. Die Evidence selbst enthält noch keine
autoritative Engineering-Aussage und stellt auch keine Persona-Interpretation
dar.

Die technische Evidence Detection ist als eigenständige Verarbeitungsstufe
realisiert. Sie erzeugt explizite Evidence-Artefakte, die mit ihrem
Source-, Processing-Run- und Attempt-Kontext verbunden bleiben.

Dadurch wird verhindert, dass die Herkunft einer Information erst aus einer
späteren LLM-Antwort rekonstruiert werden muss. Die Provenance wird bereits
vor der probabilistischen Interpretation etabliert.

Mehrere nachgelagerte Persona-Interpretationen können dadurch auf exakt
denselben Evidence-Gegenstand Bezug nehmen. Unterschiede zwischen ihren
Antworten entstehen somit nicht dadurch, dass verschiedene Personas
unterschiedliche Source-Fragmente erhalten, sondern durch ihre jeweilige
Interpretationsperspektive auf dieselbe evidenzbasierte Ausgangsinformation.

Die Trennung von Evidence Detection und Persona Interpretation ist im finalen
Realisierungsstand als eigene Responsibility Boundary implementiert und in
ADR-027 dokumentiert. Der konkrete Entwicklungsanlass dieser Entscheidung
wird an dieser Stelle nicht betrachtet, da er Bestandteil der späteren
Verifikations- und Reflexionsanalyse ist.

Source-grounded Evidence bildet damit die gemeinsame Nachweisbasis für die
weitere semantische Verarbeitung.


\subsubsection{Engineering Subject Construction}
\label{sec:engineering_subjects_impl}

Evidence allein beschreibt zunächst nur, welche Source Information für die
weitere Verarbeitung relevant ist. Für den Vergleich mehrerer
Interpretationen wird zusätzlich ein stabiler fachlicher Bezugspunkt
benötigt.

Hierzu werden aus der source-grounded Evidence
\emph{Engineering Subjects} gebildet.

Ein Engineering Subject repräsentiert den fachlichen Gegenstand, auf den sich
nachfolgende Interpretationen beziehen. Es ist damit von der Evidence selbst
zu unterscheiden: Evidence dokumentiert die beobachtete Source Information,
während das Engineering Subject den Gegenstand bezeichnet, über dessen
Bedeutung oder Einordnung unterschiedliche Interpretationen entstehen können.

Die Subject Construction verfolgt dabei das Ziel, denselben fachlichen
Gegenstand nicht künstlich durch unterschiedliche Personas oder wiederholte
Verarbeitungsausführungen zu vervielfachen. Mehrere Interpretationen sollen
stattdessen auf einem gemeinsamen Subject zusammengeführt und miteinander
vergleichbar werden.

Die Implementierung berücksichtigt deshalb nicht nur einfache
bezeichnungsartige Entitäten. Auch Engineering-Inhalte, deren Bedeutung erst
aus ihrem Source-Kontext hervorgeht, müssen als eigenständige Subjects
erhalten bleiben. Die hierfür verwendete context-preserving Subject Discovery
ist im Projekt als eigene Architektur- und Realisierungsentscheidung
dokumentiert.

Ein vereinfachtes Beispiel verdeutlicht die Trennung. Enthält eine Source
beispielsweise die Aussage

\begin{quote}
\textit{``The remote expert may control the microscope.''}
\end{quote}

so bildet der konkret referenzierte Source-Ausschnitt zunächst die
\emph{Evidence}. Als Engineering Subject kann daraus beispielsweise
\emph{Remote Expert} hervorgehen. Auf diesen gemeinsamen fachlichen
Gegenstand können anschließend unterschiedliche Persona-Interpretationen
bezogen werden. Eine Perspektive kann den Remote Expert etwa als externen
Actor einordnen, eine andere seine funktionale Verantwortung im
Remote-Control-Kontext betrachten oder eine mögliche Beziehung zu einem
System Use Case untersuchen.

Das Engineering Subject enthält damit noch keine endgültige
Modellentscheidung. Es schafft vielmehr einen stabilen fachlichen
Bezugspunkt, auf dem unterschiedliche Interpretationen verglichen,
Abweichungen sichtbar gemacht und nachfolgende Human-Review-Entscheidungen
gebunden werden können.

Mit der Bildung der Engineering Subjects ist die Source Information so weit
strukturiert, dass die eigentliche probabilistische Interpretation beginnen
kann. Der folgende Abschnitt betrachtet deshalb die Persona-basierte
Interpretation, den Vergleich der resultierenden Perspektiven sowie die
semantische Verarbeitung der Subjects.

\subsection{Persona Interpretation und semantische Verarbeitung}
\label{sec:persona_processing_impl}

Mit den source-grounded Evidence-Artefakten und den daraus gebildeten
Engineering Subjects liegt ein stabiler Bezugsrahmen für die
probabilistische Interpretation vor. Erst an dieser Stelle wird ein
Large Language Model für die fachliche Interpretation der zuvor
eingegrenzten Engineering-Information eingesetzt.

Die Implementierung verfolgt dabei nicht das Ziel, aus einem einzelnen
LLM-Aufruf unmittelbar eine autoritative Engineering-Aussage abzuleiten.
Stattdessen werden mehrere explizit konfigurierte
Interpretationsperspektiven auf denselben Source-Kontext angewendet.
Übereinstimmung, Abweichung und Instabilität werden als eigene
Verarbeitungsevidence erhalten und anschließend für den Human Engineering
Review aufbereitet.

Die semantische Verarbeitung bleibt damit von der späteren
Engineering Authority getrennt. Auch ein technisch konsistentes oder
zwischen mehreren Personas übereinstimmendes Ergebnis stellt zunächst
ausschließlich Processing Evidence dar.


\subsubsection{Source-anchored Persona Interpretation}
\label{sec:persona_calls}

Die probabilistische Interpretation erfolgt über konfigurierte
\emph{Personas}. Eine Persona beschreibt dabei keine eigenständige
Engineering Authority, sondern eine definierte Interpretationsperspektive,
aus der dieselbe Engineering-Information betrachtet wird.

Die Persona-Konfiguration legt insbesondere fest, welche fachliche
Perspektive bei der Interpretation einzunehmen ist. Der eigentliche
Source-Gegenstand wird dadurch jedoch nicht verändert. Alle für einen
gemeinsamen Analysekontext ausgeführten Personas erhalten denselben
source-gebundenen Ausgangszustand.

Technisch wird die Interpretation deshalb an eine explizite
\emph{Source Analysis Unit} gebunden. Diese referenziert den konkreten
Source-Ausschnitt und dessen Provenance und bildet den persona-unabhängigen
Analysegegenstand für die nachgelagerte Evidence Detection und Interpretation.

Die Persona selbst darf weder die Identität noch die Source-Bindung dieser
Analyseeinheit erzeugen oder verändern.

Eine solche Source Analysis Unit kann durch mehrere Personas und, sofern für
die Stabilitätsbetrachtung vorgesehen, auch durch mehrere Ausführungen
derselben Persona verarbeitet werden. Jede resultierende Interpretation
bleibt mit der konkreten Persona, dem jeweiligen Run sowie der zugrunde
liegenden Evidence verbunden.

Die erzeugten Antworten werden nicht ausschließlich als Freitext
weitergereicht. Für die weitere Verarbeitung werden relevante Aussagen
in strukturierte Interpretationsartefakte überführt. Dadurch können
verschiedene Persona-Ergebnisse auf demselben Engineering Subject
verglichen werden, ohne ihre jeweilige Herkunft zu verlieren.

Die source-gebundene Ausführung der Personas ist Bestandteil der in
ADR-026 dokumentierten Realisierung. Entscheidend für den finalen
Verarbeitungszustand ist dabei, dass semantischer Vergleich nicht auf
einem ungebundenen Pool voneinander unabhängiger LLM-Antworten erfolgt,
sondern auf Interpretationen eines gemeinsamen, explizit referenzierten
Source-Kontexts.

Die vollständigen Prompt- oder Persona-Definitionen sind für die
Architekturargumentation nicht erforderlich. Maßgeblich ist ihre
technische Rolle als kontrollierte, voneinander unterscheidbare
Interpretationsperspektiven.


\subsubsection{Consensus, Variance und Stabilität}
\label{sec:consensus_variance_impl}

Nach der Persona Interpretation werden die resultierenden Perspektiven
nicht unmittelbar zu einer gemeinsamen Aussage verdichtet. Zunächst wird
erfasst, in welchen Punkten die Interpretationen übereinstimmen und an
welchen Stellen Unterschiede bestehen.

Dabei werden drei unterschiedliche Beobachtungen voneinander getrennt:
\emph{Consensus} beschreibt die Übereinstimmung unterschiedlicher
Persona-Perspektiven, \emph{Variance} macht Unterschiede zwischen diesen
Perspektiven sichtbar und \emph{Stability} beschreibt die Reproduzierbarkeit
einer Perspektive über wiederholte Ausführungen derselben Persona.

Diese Unterscheidung ist insbesondere für wiederholte LLM-Ausführungen
relevant. Mehrere Runs derselben Persona werden nicht als mehrere
unabhängige Stimmen behandelt. Eine Persona kann für ein Engineering
Subject höchstens eine eigenständige Interpretationsperspektive beitragen.

Wiederholte Runs derselben Persona dienen stattdessen dazu, die
\emph{intra-Persona stability} zu beurteilen. Liefert dieselbe Persona bei
wiederholter Verarbeitung desselben Ausgangszustands unterschiedliche
Interpretationen, bleibt diese Abweichung als eigene Processing Evidence
sichtbar.

Damit wird verhindert, dass die Anzahl technischer LLM-Aufrufe das Gewicht
einer fachlichen Perspektive beeinflusst. Die Konsensbildung erfolgt auf
Ebene unterschiedlicher Persona-Perspektiven und nicht auf Ebene der
Anzahl ausgeführter Runs. Diese Entscheidung ist in ADR-025 und ADR-026
festgehalten.

Auch vollständige Übereinstimmung besitzt dabei keine Engineering
Authority. Ein hoher Consensus kann den nachfolgenden Review unterstützen,
ersetzt jedoch keine menschliche Entscheidung. Umgekehrt stellt Varianz
nicht automatisch einen Fehler der Verarbeitung dar. Unterschiedliche
Interpretationen können gerade bei mehrdeutiger Legacy-Information einen
relevanten Hinweis für den Human Review darstellen.


\subsubsection{Semantische Konsolidierung und source-lokale Synthese}
\label{sec:semantic_consolidation_impl}

Die einzelnen Persona-Interpretationen müssen anschließend dahingehend
geordnet werden, welche Aussagen sich auf denselben fachlichen Gegenstand
beziehen. Dazu implementiert der Proof of Concept eine semantische
Konsolidierung.

Die Konsolidierung besitzt eine begrenzte Aufgabe. Sie darf feststellen,
dass mehrere vorhandene Interpretationen denselben Engineering Subject
adressieren, sie darf jedoch keine neue Engineering-Aussage erzeugen und
keine bestehende Interpretation autoritativ auswählen.

Semantische Zusammenführung erfolgt deshalb konservativ. Eine automatische
Konsolidierung ist nur zulässig, wenn die vorhandene Vergleichsevidence
eine hinreichend eindeutige Äquivalenz unterstützt. Als unterschiedlich
oder unsicher bewertete Gegenstände bleiben getrennt und für den späteren
Human Review sichtbar.

Die in ADR-025 dokumentierte Merge Authority folgt damit einem
Fail-closed-Prinzip: Unsicherheit oder fehlende Vergleichsevidence
berechtigen nicht zu einer semantischen Zusammenführung. Das Ergebnis
einer erfolgreichen Konsolidierung bleibt ebenfalls ein abgeleitetes
Processing Artifact und keine autorisierte Engineering Information.

Die Verarbeitung erfolgt zunächst innerhalb des gemeinsamen
Source-Analysekontexts. Auf diese Weise werden die Interpretationen
verschiedener Personas lokal zu reviewbaren Engineering Subjects
zusammengeführt.

Engineering Subjects können innerhalb derselben Source jedoch an mehreren
Stellen auftreten. Nach der lokalen Konsolidierung können daher
semantisch zusammengehörige Subjects über mehrere Source Analysis Units
derselben Source hinweg synthetisiert werden. Diese
\emph{Cross-Unit Semantic Synthesis} arbeitet auf bereits lokal
konsolidierten Subjects und nicht erneut auf dem vollständigen Pool
einzelner Persona-Antworten.

Wesentlich ist die Abgrenzung zur späteren Multi-Source-Verarbeitung:
Die hier beschriebene Synthese dient der semantischen Kontinuität innerhalb
des source-lokalen Processing. Sie führt noch keine Engineering Information
unterschiedlicher registrierter Sources auf Projektebene zusammen.

Sämtliche beitragenden Evidence-, Persona- und Run-Beziehungen bleiben auch
nach der Synthese erhalten. Ein synthetisiertes Engineering Subject ersetzt
daher nicht seine Ausgangsevidence, sondern bündelt deren semantischen
Bezug für die nachfolgende Review-Grenze.

ADR-026 dokumentiert diese source-gebundene Konsolidierungs- und
Syntheselogik im finalen Realisierungszustand. Die Entstehung dieser
Entscheidung aus vorhergehenden Untersuchungsergebnissen wird erst im
Zusammenhang mit der Verifikation und Reflexion betrachtet.


\subsubsection{Semantische Referenzen und kontrollierte Klassifikation}
\label{sec:semantic_governance_impl}

Neben den unterschiedlichen Persona-Perspektiven kann die Interpretation
auf kontrollierte semantische Referenzinformation zurückgreifen. Dazu
gehören beispielsweise projekt- oder domänenspezifische Terminologie,
zulässige Klassifikationen und weitere explizit bereitgestellte
Referenzkontexte.

Diese Referenzinformation dient als semantische Leitplanke für die
Interpretation. Sie besitzt jedoch keine eigenständige Engineering
Authority und darf insbesondere nicht als Ersatz für source-grounded
Evidence verwendet werden.

Technisch werden solche semantischen Referenzen als explizit versionierte
Referenzzustände behandelt. Dazu gehören unter anderem projektbezogene
Glossary-Zustände, Terminology-Mapping-Artefakte sowie gebundene externe
Referenzressourcen. Nachgelagerte Verarbeitung kann dadurch auf einen
bestimmten semantischen Referenzstand verweisen, statt implizit einen jeweils
aktuellen Wissenszustand zu verwenden.

Auch die Gültigkeit dieser Referenzen wird kontrolliert. Fehlende, veraltete,
mehrdeutige oder inkonsistente Referenzbindungen führen nicht zu einer
automatischen semantischen Korrektur, sondern bleiben als blockierender oder
reviewbarer Zustand sichtbar.

Damit bleibt nachvollziehbar getrennt, ob eine Aussage aus der
Engineering Source stammt oder ob zusätzliche Referenzinformation lediglich
zur Einordnung ihrer Bedeutung herangezogen wurde. Eine Referenzdefinition
kann beispielsweise unterstützen, unterschiedliche Benennungen demselben
Begriffskontext zuzuordnen, sie erzeugt dadurch jedoch keinen neuen
source-seitigen Engineering-Nachweis.

Auch Klassifikationsunterschiede werden nicht durch die technische
Verarbeitung verdeckt. Erkennen mehrere Personas denselben fachlichen
Gegenstand, ordnen ihn jedoch unterschiedlich ein, kann diese Varianz auf
dem gemeinsamen Engineering Subject erhalten bleiben.

Eine semantische Harmonisierung ist daher nicht mit einer automatischen
fachlichen Entscheidung gleichzusetzen. Wo mehrere zulässige oder
widersprüchliche Interpretationen verbleiben, werden diese Informationen
für den Human Engineering Review weitergegeben.

Das Ergebnis der in diesem Abschnitt beschriebenen Verarbeitung ist somit
kein freigegebener Engineering-Inhalt, sondern ein reviewbarer
Informationszustand: source-grounded Engineering Subjects mit
nachvollziehbaren Persona-Interpretationen, Consensus-, Variance- und
Stability-Evidence sowie den zugehörigen Provenance-Beziehungen.

An dieser Stelle endet die maschinelle semantische Verarbeitung vor der
ersten fachlichen Authority Boundary. Der folgende Abschnitt beschreibt,
wie diese Informationen dem Human Engineering Review zugeführt und erst
durch eine explizite menschliche Entscheidung zu Approved Engineering
Information promoviert werden.

\subsection{Human Engineering Review und Approved Engineering Information}
\label{sec:human_review_impl}

Nach Abschluss der source-lokalen semantischen Verarbeitung liegt ein
reviewbarer Informationszustand vor. Dieser umfasst die zugrunde liegende
Evidence, die gebildeten Engineering Subjects, die zugehörigen
Persona-Interpretationen sowie Consensus-, Variance- und
Stability-Informationen.

Dieser Zustand besitzt noch keine Engineering Authority. Der Übergang von
maschinell erzeugter Processing Evidence zu autorisierter Engineering
Information erfolgt deshalb über eine explizite Human-Review-Grenze.

Die Implementierung trennt dabei den eigentlichen Review-Vorgang von der
anschließenden Promotion des freigegebenen Inhalts. Eine menschliche
Entscheidung verändert weder die ursprüngliche Source noch die bereits
erzeugten Processing Artifacts. Stattdessen entstehen zusätzliche,
versionierte Entscheidungs- und Authority-Artefakte.

Diese Trennung ist in ADR-016 als zentrale Realisierungsentscheidung
dokumentiert. Sie bildet die technische Umsetzung der in Kapitel~4
beschriebenen Human-Authority-Boundary.


\subsubsection{Subject-centered Human Review}
\label{sec:subject_review}

Der Human Engineering Review arbeitet nicht auf vollständigen
LLM-Ausgaben oder kompletten Processing Reports als einem einzigen
Freigabegegenstand. Stattdessen werden die zuvor gebildeten Engineering
Subjects als unabhängig reviewbare fachliche Einheiten verwendet.

Für jedes Review Subject wird ein konkreter Review Context bereitgestellt.
Dieser enthält die source-grounded Evidence, die zugehörigen
Persona-Interpretationen, vorhandene Consensus- und Variance-Informationen
sowie die für die Nachvollziehbarkeit erforderlichen Provenance-Bezüge.

Dadurch kann der Reviewer nicht nur das vorgeschlagene fachliche Ergebnis
betrachten, sondern auch nachvollziehen, auf welcher Source Information
dieses Ergebnis basiert und welche unterschiedlichen Interpretationen
während der Verarbeitung entstanden sind.

Die Human Review Decision wird als eigenständiges persistiertes Artefakt
behandelt. Sie überschreibt weder die ursprüngliche Source noch die
maschinell erzeugten Interpretationen. Diese bleiben als unveränderte
Processing Evidence erhalten.

Der Reviewer kann einen vorgeschlagenen Inhalt übernehmen, verändern,
zurückweisen oder für eine weitere Bearbeitung offenhalten. Werden Inhalte
geändert, entsteht ein explizit human-authored Zustand, dessen Beziehung
zu den ursprünglichen Vorschlägen erhalten bleibt.

Die Review-Entscheidung ist dabei nicht nur an eine fachliche Bezeichnung
des Subjects gebunden. Sie referenziert den konkreten, unveränderlichen
Review-Gegenstand beziehungsweise dessen Fingerprint. Dadurch wird
vermieden, dass eine bestehende Entscheidung stillschweigend auf einen
später veränderten Inhalt angewendet wird.

Das im Projekt verwendete subject-centered Review Contract konkretisiert
diese Bindung weiter. Eine Subject Decision ist an den exakten
Fingerprint der jeweiligen Review Card gebunden. Wird der Review-Gegenstand
inhaltlich verändert, muss deshalb ein neuer Entscheidungszustand
entstehen.

Die Human Review Decision stellt damit den ersten Schritt dar, an dem
Engineering Authority explizit durch eine menschliche Entscheidung
entsteht. Sie ist jedoch noch von der technischen Bereitstellung des
freigegebenen Inhalts für nachgelagerte Verarbeitungsschritte zu
unterscheiden.


\subsubsection{Promotion zu Approved Engineering Information}
\label{sec:approved_engineering_impl}

Nach Abschluss eines gültigen Human Reviews kann der freigegebene
Engineering-Inhalt in einen downstream-fähigen Authority-Zustand überführt
werden.

Die Implementierung verwendet hierfür
\emph{Approved Engineering Information}. Sie repräsentiert denjenigen
Engineering-Inhalt, der auf Grundlage einer expliziten Human Decision für
die nachfolgende Architekturableitung verwendet werden darf.

Diese Promotion erfolgt nicht unmittelbar aus einem LLM-Ergebnis und auch
nicht allein aufgrund eines hohen Persona-Consensus. Voraussetzung ist eine
gültige Human Review Decision, die an den exakt geprüften Informationszustand
gebunden ist.

Der Promotion-Schritt prüft deshalb die technische Konsistenz zwischen
Review Subject, Human Decision und dem zu promotenden Inhalt. Eine
Entscheidung, deren gebundener Inhalt nicht mehr dem aktuellen
Review-Zustand entspricht, darf nicht stillschweigend weiterverwendet werden.

Der finalisierte Review-Zustand und die daraus erzeugte Approved Engineering
Information werden als eigenständige, persistierte Artefakte erhalten.
Damit bleibt nachvollziehbar, welcher fachliche Inhalt freigegeben wurde,
durch welche Human Decision diese Freigabe erfolgte und auf welcher
source-grounded Evidence der freigegebene Zustand basiert.

Nicht jeder Bestandteil des Review Context wird dabei automatisch zu
Approved Engineering Information. Confidence-Werte, Consensus-Indikatoren,
Risiken, offene Fragen oder allgemeine Processing-Metadaten bleiben
Evidence beziehungsweise Review-Unterstützung. Nur der explizit
freigegebene Engineering-Inhalt erhält downstream Engineering Authority.

Auch Relationship-Hypothesen müssen vor einer Promotion einen expliziten
menschlichen Entscheidungszustand besitzen. Der Finalization Contract
verhindert damit, dass nicht entschiedene ausgehende
Relationship-Hypothesen implizit in einen autorisierten Zustand übergehen.

Approved Engineering Information ist somit nicht als nachträglich
„bereinigter“ LLM-Output zu verstehen. Sie bildet einen neuen
Authority-Zustand, dessen Inhalt aus einer nachvollziehbaren Kette von
Source Evidence, Processing, Human Review und expliziter Promotion
hervorgeht.

Die ursprünglichen Source- und Processing-Artefakte bleiben dabei erhalten
und werden nicht durch den autorisierten Zustand ersetzt. Dadurch kann auch
nach der Promotion weiterhin zwischen Source Authority, Processing Evidence,
Human Decision und Approved Engineering Information unterschieden werden.

Mit der Approved Engineering Information endet die source-lokale
Engineering-Authority-Bildung. Erst im nächsten Verarbeitungsschritt wird
geprüft, welche dieser autorisierten Informationen im projektweiten Kontext
für die weitere Modellableitung zugelassen werden.

\subsection{Project Fit und Multi-Source-Verarbeitung}
\label{sec:multisource_impl}

Bis zur Approved Engineering Information erfolgt die Verarbeitung weiterhin
source-lokal. Auch wenn mehrere Engineering Sources demselben Projekt
zugeordnet sind, werden ihre Processing Evidence, Human Decisions und
autorisierten Engineering-Inhalte nicht während der source-lokalen
Verarbeitung zusammengeführt.

Für die projektweite Modellableitung muss deshalb zunächst entschieden werden,
welche source-lokal autorisierte Engineering Information überhaupt in den
gemeinsamen Projektkontext eingehen darf. Der Proof of Concept realisiert
hierfür eine explizite Project-Fit-Grenze.

Der final akzeptierte Multi-Source-Pfad erweitert damit die bestehende
Single-Source-Verarbeitung, ohne deren Authority- und Provenance-Strukturen
aufzulösen. Jede Source behält ihre eigene Verarbeitungshistorie und ihre
eigene Approved Engineering Information. Erst nach erfolgreicher
Project-Fit-Prüfung werden diese getrennten source-lokalen Authority-Zustände
für die gemeinsame Architekturableitung bereitgestellt.


\subsubsection{Project Fit als projektweite Zulassungsgrenze}
\label{sec:project_fit_impl}

Die Registrierung einer Source innerhalb eines Projekts reicht nicht aus, um
ihre Engineering Information automatisch für die projektweite
Modellableitung zuzulassen. Eine registrierte Source kann beispielsweise
lediglich Kontextinformation enthalten, irrtümlich einem Projekt zugeordnet
worden sein oder hinsichtlich ihrer Zugehörigkeit zum betrachteten
Engineering-Gegenstand nicht eindeutig sein.

Der Proof of Concept führt deshalb nach der source-lokalen Verarbeitung eine
eigene \emph{Project Fit Assessment}-Stufe ein. Sie beantwortet die begrenzte
Frage, ob eine Source plausibel zum Engineering-Kontext des betrachteten
Projekts gehört.

Die Bewertung kann hierzu den registrierten Project Context, verfügbare
Kontextquellen und bereits vorhandene projektbezogene Information
berücksichtigen. Die verwendeten Kontextinformationen erhalten dadurch jedoch
keine Engineering Authority. Sie unterstützen ausschließlich die
Zugehörigkeitsbewertung.

Das Project Fit Assessment erzeugt einen expliziten, persistierten
Bewertungszustand. Die maschinelle Bewertung unterscheidet dabei zwischen

\begin{quote}
\texttt{plausible\_in\_scope},\\
\texttt{uncertain} und\\
\texttt{likely\_out\_of\_scope}.
\end{quote}

Nur ein ausreichend bestätigter Project-Fit-Zustand erlaubt die Übergabe der
betroffenen source-lokalen Engineering Information an die projektweite
Modellableitung. Unsichere, fehlende oder nicht mehr zum gebundenen
Artefaktzustand passende Project-Fit-Evidence darf diesen Übergang nicht
stillschweigend autorisieren.

Die Prüfung erfolgt deshalb nicht lediglich anhand einer Source ID oder eines
Dateinamens. Der Project-Fit-Zustand wird an den konkreten Project-, Source-,
Projection-, Processing-Run- und Attempt-Kontext gebunden. Dadurch bleibt
nachvollziehbar, für welchen exakten Verarbeitungszustand die
Zugehörigkeitsentscheidung getroffen wurde.

Die Project-Fit-Grenze besitzt dabei selbst keine Engineering Authority über
den Inhalt der Source. Eine positive Bewertung erklärt weder eine
LLM-Interpretation für fachlich korrekt noch verändert sie die zuvor
getroffene Human Review Decision. Sie entscheidet ausschließlich über die
Zulässigkeit der bereits autorisierten source-lokalen Information für den
nachfolgenden projektweiten Verarbeitungspfad.

Damit bleiben zwei unterschiedliche Entscheidungen getrennt:
Der source-lokale Human Engineering Review entscheidet, \emph{welcher
Engineering-Inhalt autorisiert ist}; Project Fit entscheidet anschließend,
\emph{ob dieser autorisierte Inhalt in den betrachteten Projektkontext
eingebracht werden darf}.


\subsubsection{Multi-Source Handoff und getrennte AEI-Konsumtion}
\label{sec:multisource_handoff_impl}

Nach erfolgreicher Project-Fit-Prüfung wird der zulässige source-lokale
Authority-Zustand über einen expliziten technischen Handoff an die
Modellableitung übergeben.

Der dafür realisierte \emph{ProjectFitPhaseHHandoff} bindet die Project-Fit-
Entscheidung an die zugehörigen Approved Inputs beziehungsweise die
source-lokale Approved Engineering Information. Zusätzlich werden die
relevanten Artefakt-Fingerprints erhalten. Der nachgelagerte
Verarbeitungsschritt kann dadurch prüfen, dass genau der freigegebene und
für das Projekt zugelassene Informationszustand konsumiert wird.

Bei mehreren Sources wird dabei keine neue synthetische
Approved Engineering Information erzeugt. Jede source-lokale AEI bleibt als
eigenständiger autorisierter Informationszustand mit ihrer ursprünglichen
Source- und Processing-Provenance erhalten.

Die projektweite Verarbeitung konsumiert diese AEIs deshalb getrennt. Ihre
gemeinsame Verwendung für die Architekturableitung bedeutet nicht, dass ihre
Engineering-Aussagen zuvor zu einem einzigen Authority-Artefakt verschmolzen
werden.

Auch die technische Identität bleibt source- beziehungsweise run-spezifisch.
An projektweiten Aggregationsgrenzen wird ein Processing Artifact über

\[
(\texttt{processing\_run\_id},
 \texttt{artifact\_type},
 \texttt{artifact\_id})
\]

identifiziert. Gleichlautende lokale Artifact IDs aus unterschiedlichen
Processing Runs können dadurch ohne Identitätskollision nebeneinander
bestehen.

Für Engineering Subjects wird die Source-Zugehörigkeit ebenfalls erhalten.
Damit kann derselbe oder ein ähnlich bezeichneter fachlicher Gegenstand aus
unterschiedlichen Sources weiterhin auf seine jeweilige source-lokale
Authority und Evidence zurückgeführt werden.

Der Multi-Source-Pfad führt folglich nicht über eine automatische
Cross-Source-Wahrheitsentscheidung. Weder die Anzahl übereinstimmender Sources
noch Source-Reihenfolge, Source-Alter oder maschinell erkannte semantische
Ähnlichkeit dürfen eine bestehende Engineering Authority automatisch
überschreiben.

Die weitergehende concern-zentrierte Cross-Source-Reconciliation wurde im
Projekt als zusätzlicher Forschungs- und Evolutionspfad untersucht. Sie ist
jedoch nicht Bestandteil der verpflichtenden Verarbeitungskette des finalen
Thesis-Prototyps. Für den hier beschriebenen Realisierungsstand endet die
Multi-Source-Zulassung daher mit dem exakten Project-Fit-Handoff und der
getrennten Bereitstellung der zugelassenen source-lokalen Approved Engineering
Information.

Auf dieser Grundlage kann im nächsten Schritt aus mehreren zugelassenen
Engineering-Informationsbeiträgen genau ein projektweites
\emph{Model Candidate Set} abgeleitet werden. Die dabei entstehenden
Modellvorschläge besitzen wiederum noch keine Model Authority und werden vor
der Assemblierung einer eigenen Human-Model-Placement-Grenze zugeführt.

\subsection{Model Candidate Derivation und Human Model Placement}
\label{sec:model_candidate_impl}

Nach erfolgreichem Project Fit stehen die für die Modellableitung zugelassenen
source-lokalen Approved Engineering Information-Artefakte zur Verfügung.
Damit beginnt der Übergang von autorisierter Engineering Information zu einer
konkreten Modellrepräsentation.

Dieser Übergang erfolgt nicht unmittelbar durch die Erzeugung von
SysML-v2-Elementen. Der Proof of Concept führt zunächst eine eigenständige
Model-Candidate-Ebene ein. Sie dient dazu, aus der autorisierten
Engineering Information Vorschläge für Modellelemente und Beziehungen
abzuleiten, ohne bereits Model Authority zu erzeugen.

Zwischen der fachlichen Ableitung eines Model Candidates und seiner
strukturellen Einordnung in das Zielmodell besteht eine weitere
Authority Boundary. Die konkrete Platzierung eines fachlichen Gegenstands
innerhalb der vorgesehenen Modellstruktur wird deshalb separat betrachtet und
durch einen Human Model Placement Review autorisiert.

Damit werden zunächst drei Zustände voneinander getrennt:
autorisierte Engineering Information, daraus abgeleitete Modellvorschläge und
menschlich autorisierte Model Placement Information. Erst auf Grundlage dieser
Zustände kann anschließend ein konkreter interner Modellvorschlag assembliert
und einer weiteren menschlichen Prüfung zugeführt werden.


\subsubsection{Ableitung des projektweiten Model Candidate Set}
\label{sec:model_candidate_set}

Die Modellableitung konsumiert ausschließlich Engineering Information, die den
vorhergehenden Human-Authority- und Project-Fit-Grenzen erfolgreich
durchlaufen hat. Draft-Zustände, nicht freigegebene Persona-Interpretationen
oder lediglich maschinelle Consensus-Bewertungen stellen keine gültigen
Eingaben für diesen Verarbeitungsschritt dar.

Aus den zugelassenen source-lokalen Informationsbeiträgen wird ein gemeinsamer
projektweiter Modellvorschlag aufgebaut. Die technische Repräsentation hierfür
bildet das \emph{Model Candidate Set}.

Ein Model Candidate Set fasst die für einen konkreten Ableitungszustand
erzeugten Element- und Relationship-Candidates zusammen und bindet sie an die
exakten zugrunde liegenden Engineering-Informationen. Bei einer
Multi-Source-Verarbeitung werden dabei weiterhin die separaten
source-lokalen Provenance-Beziehungen erhalten. Die Kandidatenbildung erzeugt
keine synthetische Approved Engineering Information.

Element Candidates repräsentieren vorgeschlagene Modellgegenstände. Sie
enthalten unter anderem einen fachlichen Bezug, eine vorgeschlagene
Bezeichnung und Beschreibung sowie die für die spätere strukturelle
Einordnung benötigte Klassifikations- und Provenance-Information.

Relationship Candidates werden als eigenständige Artefakte behandelt.
Beziehungen werden damit nicht lediglich implizit als Eigenschaften eines
Element Candidates gespeichert. Source- und Target-Bezüge sowie die
zugrunde liegende semantische Bedeutung bleiben explizit nachvollziehbar.

Ein Model Candidate stellt dabei keine neue Engineering Authority dar.
Die bereits autorisierte Engineering Information bleibt die fachliche
Ausgangsautorität. Der Candidate beschreibt lediglich eine mögliche
Überführung dieser Information in eine modellierbare Struktur.

Auch alternative Modellierungsoptionen können als getrennte Candidates
erhalten bleiben. Mehrdeutigkeit wird dadurch nicht bereits während der
Ableitung durch eine technisch bevorzugte Variante verdeckt.

Die Einführung dieser eigenständigen Modellvorschlagsebene ist in ADR-018
dokumentiert. Sie trennt die fachlich autorisierte Eingangsinformation von
ihrer späteren strukturellen und repräsentationsspezifischen Umsetzung.

Das Model Candidate Set bildet damit einen reproduzierbaren
Ableitungszustand. Nachgelagerte Verarbeitung bezieht sich auf einen explizit
ausgewählten Candidate-Zustand und nicht implizit auf den jeweils zuletzt
erzeugten Modellvorschlag.


\subsubsection{Human Model Placement und Approved Model Placement Set}
\label{sec:model_placement_impl}

Die Ableitung eines Model Candidates beantwortet noch nicht abschließend, an
welcher Stelle der vorgesehenen Modellstruktur der jeweilige fachliche
Gegenstand eingeordnet werden soll.

Diese Frage wird in einer eigenständigen
\emph{Model Placement}-Stufe behandelt. Sie bestimmt beispielsweise den
betroffenen Modellbereich, die vorgesehene Elementklasse sowie die
strukturelle Position innerhalb des verwendeten Model Structure Profile.

Für die Platzierungsunterstützung können sowohl deterministische
Profilinformationen als auch LLM-basierte Modellierungsperspektiven verwendet
werden. Im LLM-unterstützten Pfad werden mehrere speziell auf
Modellierungsentscheidungen ausgerichtete Personas eingesetzt. Alle
Perspektiven erhalten denselben autorisierten Engineering-Gegenstand und
dieselben zulässigen Placement-Optionen.

Auch hier erzeugt Persona Agreement keine Model Authority. Übereinstimmende
Placement-Vorschläge können die Review-Entscheidung unterstützen, während
abweichende oder nicht eindeutige Vorschläge als sichtbare Varianz erhalten
bleiben.

Die resultierenden Placement-Vorschläge werden deshalb in einem
\emph{Human Model Placement Review} zusammengeführt. Der Reviewer entscheidet
für jeden relevanten Engineering-Gegenstand, welche der profilkonformen
Platzierungsoptionen autorisiert wird.

Die menschliche Placement Decision wird an den exakten Informations-,
Candidate- und Vorschlagszustand gebunden. Dadurch kann eine einmal getroffene
Platzierungsentscheidung nicht stillschweigend auf später veränderte Candidates
oder Profilzustände übertragen werden.

Eine vollständig aufgelöste Menge solcher Entscheidungen bildet das
\emph{Approved Model Placement Set}. Dieses Artefakt repräsentiert die
autorisierte strukturelle Einordnung der für die Modellassemblierung
vorgesehenen Engineering-Information.

Die Trennung von Placement und Assembly ist in ADR-029 festgelegt.
Entscheidend ist dabei, dass strukturelle Mehrdeutigkeit nicht durch die
Assemblierungslogik selbst aufgelöst werden darf. Model Assembly erhält
bereits menschlich autorisierte Platzierungsinformation und muss deshalb
keine eigenständige Entscheidung darüber treffen, wo ein Engineering Subject
im Modell einzuordnen ist.

Akzeptierte semantische Beziehungen aus der Approved Engineering Information
werden dabei nicht allein deshalb erneut interpretiert, weil ihre beteiligten
Elemente strukturell platziert werden. Ihre fachliche Semantik bleibt erhalten
und wird im nachfolgenden Assembly-Schritt mit den autorisierten
Elementplatzierungen verbunden.

Mit dem Approved Model Placement Set liegt damit ein autorisierter
Strukturierungszustand vor. Dieser ist jedoch noch nicht mit dem finalen
generierbaren Internal Engineering Model gleichzusetzen. Zunächst wird aus den
autorisierten Placements ein konkreter Assembly-Zustand erzeugt und nochmals
als zusammenhängender Modellvorschlag geprüft.

\subsection{Internal Engineering Model, Target-Model Formulation und SysML-v2-Generierung}
\label{sec:internal_model_generation}

Nach Abschluss des Human Model Placement liegt eine autorisierte strukturelle
Einordnung der Engineering-Information vor. Der final implementierte
Verarbeitungspfad führt von diesem Zustand jedoch nicht unmittelbar zur
SysML-v2-Generierung.

Die Umsetzung unterscheidet vielmehr zwischen der deterministischen
Assemblierung eines ersten internen Modellzustands, dessen menschlicher
Freigabe, einer nachgelagerten Target-Model Formulation sowie einer weiteren
Human-Authority-Grenze für die resultierende Modellqualität.

Erst der daraus entstehende human-authorized Successor des
Internal Engineering Model bildet den normativen Eingang für die
deterministische SysML-v2-Generierung.

Damit wird zwischen Engineering Meaning, struktureller Platzierung,
Target-Model Formulation und textueller Serialisierung unterschieden.


\subsubsection{Model Assembly Review und Base Internal Engineering Model}
\label{sec:base_iem_assembly}

Auf Grundlage des Approved Model Placement Set wird zunächst ein konkreter
Model Assembly Draft erzeugt. Die Assemblierungslogik ist deterministisch und
darf keine zusätzliche Engineering-Semantik erzeugen.

Die autorisierten Engineering-Gegenstände werden entsprechend ihrer
freigegebenen Platzierungsinformation in die ausgewählte Framework-Struktur
eingebettet. Auch akzeptierte Relationships werden mit ihren bereits
autorisierten semantischen Eigenschaften übernommen.

Interne Modellelemente und Beziehungen erhalten dabei eigenständige technische
Identitäten. Candidate Identity, Approved Engineering Information und interne
Model Identity bleiben getrennt, werden jedoch über explizite
Traceability-Beziehungen miteinander verbunden.

Die Assembly darf weder eine fehlende Placement Decision ersetzen noch eine
strukturelle Mehrdeutigkeit selbständig auflösen. Ist eine autorisierte
Platzierung nicht eindeutig materialisierbar, bleibt der Verarbeitungspfad
blockiert.

Der erzeugte Assembly-Zustand wird anschließend in einem eigenständigen
\emph{Model Assembly Review} als zusammenhängender Modellvorschlag geprüft.
Die Human Decision wird an den konkreten Assembly Draft und dessen
Fingerprint gebunden.

Erst nach dieser menschlichen Freigabe wird der Assembly-Zustand als
authority-backed \emph{Base Internal Engineering Model} materialisiert.

Das Base IEM repräsentiert damit einen reproduzierbaren und menschlich
freigegebenen internen Modellzustand. Es ist representation-neutral in dem
Sinne, dass seine Engineering-Semantik noch nicht mit einer konkreten
SysML-v2-Textserialisierung gleichgesetzt wird.

Der freigegebene Base-IEM-Zustand bleibt immutable. Spätere Änderungen erzeugen
keine stille Mutation dieses Snapshots, sondern einen neuen nachgelagerten
Modellzustand.

Der Proof of Concept besitzt hierfür eine explizite Final Assembly Decision.
Diese menschliche Entscheidung ist Bestandteil der Authority-Bindung des
materialisierten Internal Engineering Model.

\subsubsection{Target-Model Formulation, Model Refinement und Human Model Quality Review}
\label{sec:target_model_formulation_impl}

Das freigegebene Base Internal Engineering Model enthält die autorisierte
Engineering-Semantik und ihre strukturelle Platzierung. Für eine geeignete
Repräsentation im vorgesehenen Target Model können jedoch zusätzliche
Entscheidungen darüber erforderlich sein, wie dieser Inhalt in der
unterstützten Zielmodellstruktur repräsentiert werden soll.

Der Proof of Concept trennt deshalb
\emph{Engineering Meaning},
\emph{Target-Model Representation} und
\emph{Target-Model Formulation} voneinander.

Für die Target-Model Formulation werden zunächst explizite
Repräsentationsoptionen für den gebundenen Base-IEM-Zustand gebildet. Diese
können beispielsweise festlegen, welches unterstützte Target-Model-Pattern oder
welches SysML-v2-Konstrukt für einen Engineering-Gegenstand verwendet werden
soll. Die Vorschläge besitzen zunächst keine Model Authority.

Die Auswahl erfolgt über eine eigenständige menschliche
\emph{Target-Model Formulation Decision}. Die Entscheidungen werden an den
konkreten Review-Gegenstand, den ausgewählten Candidate sowie die verwendeten
Target-Model- und Target-Notation-Referenzen gebunden. Erst eine vollständige
Menge solcher Entscheidungen bildet die
\emph{Target-Model Formulation Authority} für den betrachteten Modellzustand.

Damit wird bereits vor einer weiteren probabilistischen Verarbeitung
menschlich festgelegt, welche Target-Model-Repräsentation für den jeweiligen
Engineering-Gegenstand zulässig ist.

Auf Grundlage dieses autorisierten Zustands führt der Proof of Concept
anschließend eine begrenzte \emph{Model-Quality-Refinement}-Stufe aus. Dieser
Schritt kann LLM-basiert erfolgen und dient insbesondere dazu, Namen und
Beschreibungen für die bereits autorisierte Zielrepräsentation zu verbessern.

Die Refinement-Stufe besitzt keine Authority, die vorgelagerte
Engineering-Bedeutung oder Target-Model-Entscheidung zu verändern. Sie arbeitet
innerhalb der bereits festgelegten Repräsentationsgrenzen und darf keine neue
Engineering-Semantik hinzufügen.

Die erzeugten Refinement Proposals werden anschließend in einem
\emph{Human Model Quality Review} geprüft. Für jedes betroffene Modellelement
wird eine explizite menschliche Entscheidung erzeugt. Ein Vorschlag kann
übernommen, durch eine human-authored Formulierung überschrieben oder
zurückgewiesen werden.

Eine vollständige Menge freigegebener Entscheidungen bildet die
\emph{Model Quality Authority}. Diese zweite Authority-Stufe ist an den exakten
Refinement-Zustand und dessen Fingerprints gebunden.

Der implementierte Verarbeitungspfad trennt damit zwei unterschiedliche
menschliche Entscheidungen: Die Target-Model Formulation Authority legt fest,
\emph{wie ein autorisierter Engineering-Gegenstand grundsätzlich im Zielmodell
repräsentiert werden soll}. Die nachgelagerte Model Quality Authority
entscheidet, \emph{welche konkrete Formulierung innerhalb dieser bereits
autorisierten Repräsentation verwendet werden darf}.

Erst wenn beide Authority-Zustände vollständig vorliegen, kann daraus ein
neuer Internal-Engineering-Model-Zustand materialisiert werden.

\subsubsection{Human-authorized Successor Internal Engineering Model}
\label{sec:successor_iem_impl}

Erst nach erfolgreichem Human Model Quality Review wird der akzeptierte
Target-Model-Zustand in einen neuen Internal-Engineering-Model-Snapshot
überführt.

Dieser \emph{Successor IEM} ersetzt den Base-IEM-Zustand nicht destruktiv.
Beide Modellzustände bleiben als getrennte immutable Snapshots erhalten und
sind über ihre Provenance- und Authority-Beziehungen miteinander verbunden.

Die Materialisierung des Successor IEM benötigt sowohl die freigegebene
Target-Model Formulation als auch die zugehörige Human Model Quality
Authority. Damit kann kein ausschließlich maschinell erzeugter
Formulation-Vorschlag selbständig zum neuen generierbaren Modellzustand werden.

Der Successor übernimmt die weiterhin gültige Engineering-Semantik des
Base IEM und ergänzt beziehungsweise konkretisiert ausschließlich diejenigen
Repräsentationsaspekte, die durch die nachgelagerte Human-Authority-Kette
freigegeben wurden.

Auch Relationship-Repräsentationen bleiben an ihre akzeptierte Engineering-
Semantik gebunden. Ändert eine Target-Model Formulation die effektive
Repräsentation eines beteiligten Modellelements, dürfen daraus entstehende
Kompatibilitätsfragen nicht durch eine implizite Generatorentscheidung
aufgelöst werden.

Der resultierende Successor IEM bildet damit den exakten autorisierten
Modellzustand, der anschließend an die SysML-v2-Generierung übergeben wird.

Die nachfolgende Generation muss daher weder Engineering Meaning noch
Target-Model Formulation erneut bestimmen. Ihre Aufgabe ist ausschließlich die
deterministische Repräsentation des bereits autorisierten Successor-IEM-
Zustands.


\subsubsection{Deterministische SysML-v2-Generierung}
\label{sec:sysml_generation_impl}

Die SysML-v2-Generierung konsumiert einen explizit ausgewählten und
human-authorized Internal-Engineering-Model-Snapshot. Im finalen
Verarbeitungspfad ist dies der zuvor erzeugte Successor IEM.

Eine implizite Auswahl des jeweils neuesten Modellzustands ist nicht
vorgesehen. Der konkret zu serialisierende IEM-Zustand wird über seine
technische Identität und seinen Fingerprint adressiert.

Die Generation ist als deterministische Serialisierung implementiert.
Large Language Models oder andere probabilistische Verfahren besitzen in
dieser Stufe keine Rolle für die Auswahl von Engineering-Semantik,
Modellelementen oder Beziehungen.

Die Abbildung des IEM auf SysML v2 erfolgt über versionierte
Generierungsregeln. Dabei werden drei unterschiedliche Regelbereiche
auseinandergehalten: die zulässige Zielnotation, die semantische Abbildung
interner Modellkonzepte auf SysML-v2-Konstrukte sowie die Organisation der
erzeugten Artefakte.

Die verwendete Zielnotation orientiert sich an der normativen
SysML-v2-Spezifikation \cite{SysMLv2}. Der Proof of Concept unterstützt dabei
bewusst nur einen begrenzten, explizit definierten Ausschnitt der verfügbaren
Sprachkonstrukte.

Vor der eigentlichen Erzeugung wird geprüft, ob für alle im IEM enthaltenen
und zur Materialisierung vorgesehenen Engineering-Semantiken eine eindeutige
und unterstützte Mapping-Regel vorliegt. Fehlt eine erforderliche Abbildung
oder ist eine Zuordnung mehrdeutig, wird die Generation blockiert.

Der Generator verwendet insbesondere keine generischen Fallbacks, um
unbekannte Engineering-Semantik dennoch syntaktisch darstellbar zu machen.
Eine nicht unterstützte Beziehung darf beispielsweise nicht allein deshalb
als Dependency ausgegeben werden, weil dieses Konstrukt syntaktisch verfügbar
ist.

Die Zielnotation bleibt damit an die zuvor autorisierte Engineering-Semantik
gebunden. Eine erfolgreiche Generation darf keine Bedeutung hinzufügen,
entfernen oder stillschweigend verändern.

Auch die erzeugten Symbole werden deterministisch aus stabilen internen
Identitäten abgeleitet und nicht allein aus fachlichen Anzeigenamen gebildet.
Dadurch bleiben technische Referenzen eindeutig, selbst wenn mehrere
Engineering-Elemente ähnliche oder identische Bezeichnungen besitzen.

Source-seitige Namen, Beschreibungen und weitere Textinhalte werden bei der
Serialisierung als Daten behandelt. Sie werden nicht ungeprüft als freie
SysML-v2-Syntax in das Zielartefakt übernommen.

Der erzeugte Output wird als \emph{Generated SysML Artifact Set}
repräsentiert. Die Architektur unterstützt grundsätzlich mehrere
Ausgabeeinheiten, während der im Proof of Concept verwendete
Realisierungsstand eine kompakte textuelle SysML-v2-Repräsentation erzeugt.

Listing~\ref{lst:generated_sysml_example} zeigt exemplarisch einen Ausschnitt
eines durch den Proof of Concept erzeugten SysML-v2-Artefakts. Sichtbar werden
dabei unterschiedliche, deterministisch aus dem autorisierten internen
Modellzustand erzeugte SysML-v2-Konstrukte. Die technischen
Elementbezeichner werden aus stabilen internen Modellidentitäten abgeleitet,
während fachlicher Name und Beschreibung als dokumentierende Information
erhalten bleiben.

\begin{lstlisting}[
    caption={Exemplarischer Ausschnitt eines durch den Proof of Concept
    erzeugten SysML-v2-Artefakts},
    label={lst:generated_sysml_example}
]
requirement IME_000010 {
    doc /* Engineering name: Operator Permission for Temporary Control
    Description:
    Temporary microscope control by the remote expert shall require
    operator permission. */
}

action IME_000009 {
    doc /* Engineering name: Take Temporary Microscope Control
    Description:
    The remote expert may take temporary control of the microscope
    when permitted by the operator. */
}
\end{lstlisting}

Die Generation bleibt damit unabhängig von der nachfolgenden
Artefaktvalidierung. Ein syntaktisch erzeugtes SysML-v2-Artefakt gilt nicht
allein aufgrund erfolgreicher Serialisierung als validiert oder
veröffentlichungsfähig.

ADR-021 dokumentiert die Trennung zwischen autorisiertem IEM,
deterministischer Serialisierung und nachgelagerter Validierung als
eigenständige Generierungsarchitektur.

Mit Abschluss dieses Schritts liegt erstmals eine konkrete
SysML-v2-Repräsentation des final autorisierten internen Modellzustands vor.
Ob dieses Artefakt die technischen und architektonischen
Validierungsbedingungen erfüllt und anschließend menschlich freigegeben
werden darf, wird im folgenden Abschnitt betrachtet.

\subsection{Generated Artifact Validation, Final Human Review und Publication}
\label{sec:generated_validation_publication}

Mit Abschluss der deterministischen SysML-v2-Generierung liegt ein konkretes
\emph{Generated SysML Artifact Set} vor. Dieses Artefakt repräsentiert den
zuvor autorisierten Internal-Engineering-Model-Zustand in der unterstützten
Zielnotation.

Eine erfolgreiche Generierung bedeutet jedoch noch nicht, dass das erzeugte
Artefakt technisch valide oder zur Veröffentlichung freigegeben ist. Der
Proof of Concept trennt deshalb drei nachgelagerte Verantwortlichkeiten:
die technische Validierung des erzeugten Artefakts, die abschließende
menschliche Prüfung des vollständigen Modellzustands und die kontrollierte
Publication.

Damit bleibt auch am Ende der Verarbeitung zwischen technischer
Validität und Engineering beziehungsweise Release Authority unterschieden.


\subsubsection{Interne und externe SysML-v2-Validierung}
\label{sec:generated_artifact_validation}

Die Validierung arbeitet auf einem explizit ausgewählten und unveränderlichen
Generated SysML Artifact Set. Sie darf weder die erzeugte Notation verändern
noch eine erneute Engineering-Interpretation durchführen.

Die Implementierung verwendet zwei komplementäre Validierungsebenen.

Die erste Ebene besteht aus deterministischen internen Prüfungen innerhalb
des Turing Generators. Sie prüfen insbesondere die Integrität des
Generated Artifact Set, die Übereinstimmung mit den gebundenen
Generierungsprofilen, die Struktur der erzeugten Einheiten, Relationship-
und Endpoint-Konsistenz sowie die Vollständigkeit der übertragenen
Traceability.

Die zweite Ebene bildet eine externe SysML-v2-Kompatibilitätsprüfung.
Hierfür wird das erzeugte Artefakt über einen dedizierten Validator-Adapter
an den SYSIDE Modeler übergeben. Die externe Prüfung dient dazu, die
Kompatibilität des vollständigen erzeugten Artefakts mit der verwendeten
SysML-v2-Toolumgebung zu überprüfen.

Interne und externe Validierung erfüllen dabei unterschiedliche Aufgaben.
Die internen Checks prüfen Turing-spezifische Verträge und
Artefaktinvarianten. Die externe Prüfung bestätigt dagegen die
Verarbeitbarkeit der erzeugten Notation in der ausgewählten
SysML-v2-Umgebung. Keine der beiden Ebenen ersetzt die jeweils andere.

Die Validierung ist rein beobachtend. Ein fehlerhaftes Artefakt wird nicht
automatisch korrigiert, um einen erfolgreichen Prüfstatus zu erreichen.
Festgestellte Abweichungen werden stattdessen als explizite
Validation Findings erhalten.

Das resultierende \emph{SysMLValidationResult} unterscheidet zwischen den
Zuständen

\begin{quote}
\texttt{valid},\\
\texttt{invalid} und\\
\texttt{incomplete}.
\end{quote}

Der Zustand \texttt{incomplete} wird verwendet, wenn eine erforderliche
Validierungsumgebung oder ein notwendiger Prüfkontext nicht ausreichend
verfügbar ist. Dadurch wird insbesondere verhindert, dass eine nicht
durchgeführte externe Prüfung fälschlich als erfolgreiche Validierung
interpretiert wird.

Nur ein vollständiges Ergebnis mit

\begin{quote}
\texttt{validation\_status = valid}
\end{quote}

und einem bestandenen Publication Gate ist für den nachfolgenden
Veröffentlichungspfad zulässig.

Auch das Validation Result wird an den exakten Generated-Artifact-Zustand,
die verwendeten Validierungsprofile und die konkrete Validator-Umgebung
gebunden. Eine spätere Änderung des generierten Inhalts führt daher zu
einem neuen Validierungsgegenstand.

Die in ADR-022 dokumentierte Validation Architecture trennt damit
Generation und Validierung bewusst voneinander und realisiert die
technische Freigabegrenze vor dem abschließenden Human Review.


\subsubsection{Final Human Model Review}
\label{sec:final_model_review_impl}

Auch ein technisch valides Generated SysML Artifact Set besitzt noch keine
abschließende Publication Authority.

Der Proof of Concept sieht deshalb nach der technischen Validierung einen
\emph{Final Human Model Review} vor. Dieser Review betrachtet erstmals den
vollständig generierten Modellzustand einschließlich seiner konkreten
SysML-v2-Repräsentation als zusammenhängendes Ergebnis.

Die Review-Oberfläche kann dabei sowohl die textuelle SysML-v2-Darstellung
als auch modellorientierte Projektionen, Beziehungen,
Validation Findings und die zugehörige Traceability bereitstellen. Der
Reviewer kann dadurch den resultierenden Modellzustand nicht nur auf
technischer, sondern auch auf fachlicher und struktureller Ebene betrachten.

Der Final Human Model Review ersetzt keine der vorgelagerten
Authority Boundaries. Source-lokaler Human Engineering Review und Human
Model Placement bleiben weiterhin die Autoritätsgrenzen für Engineering
Meaning und strukturelle Platzierung.

Der abschließende Review besitzt eine andere Aufgabe: Er bewertet den
vollständigen generierten Modellzustand vor seiner Veröffentlichung.

Eine Final Model Review Revision ist an das exakte Generated SysML Artifact
Set und das dazugehörige SysMLValidationResult gebunden. Wird der
Generierungs- oder Validierungszustand geändert, entsteht ein neuer
Review-Gegenstand.

Änderungswünsche aus diesem Review werden nicht durch direkte Mutation des
veröffentlichten oder internen Modellzustands umgesetzt. Stattdessen werden
sie als explizite Change Requests an die jeweils zuständige
Verarbeitungs- oder Authority Boundary zurückgegeben.

Betrifft eine Änderung beispielsweise die Engineering-Semantik, muss sie
über den entsprechenden Upstream-Review-Pfad erneut autorisiert werden.
Betrifft sie ausschließlich die textuelle Repräsentation, ist dagegen die
Generierungslogik beziehungsweise das zugehörige Profil zuständig.

Nach einer materiellen Änderung wird die nachgelagerte Kette erneut
durchlaufen:

\begin{quote}
autorisierter neuer Modellzustand\\
$\rightarrow$ neue Generation\\
$\rightarrow$ neue Validierung\\
$\rightarrow$ neue Final Model Review Revision.
\end{quote}

Eine frühere Validierung oder Human Approval wird damit nicht auf veränderte
Artefakte übertragen.

Der Final Human Review endet mit einer expliziten
\emph{Final Model Review Decision}. Nur eine Entscheidung, die den exakten
Review-, Artifact- und Validation-Zustand für die Veröffentlichung freigibt,
erzeugt die erforderliche Human Release Authority.

Diese Trennung ist in ADR-023 dokumentiert. Insbesondere kann eine
erfolgreiche SYSIDE-Validierung die abschließende menschliche
Freigabeentscheidung nicht ersetzen.


\subsubsection{Kontrollierte Publication}
\label{sec:controlled_publication_impl}

Die Publication bildet den letzten Schritt der prototypisch realisierten
Informationsverarbeitung.

Sie darf ausschließlich erfolgen, wenn sowohl die technische
Validierungsgrenze als auch die abschließende Human-Release-Grenze
erfolgreich durchlaufen wurden.

Vor der Publication werden deshalb die Bindungen zwischen
Generated SysML Artifact Set, SysMLValidationResult und Final Model Review
Decision erneut geprüft. Insbesondere müssen alle drei Artefakte demselben
Projekt- und Modellzustand zugeordnet sein und über die erwarteten
Fingerprints aufeinander verweisen.

Eine veraltete Human Decision darf keinen später veränderten Output
autorisieren. Ebenso darf ein Validation Result nicht für ein anderes oder
zwischenzeitlich neu generiertes Artifact Set verwendet werden.

Erst nach erfolgreicher Prüfung dieser Bedingungen wird ein
\emph{Published Output Package} erzeugt.

Nicht freigegebene Generated Artifacts werden weiterhin als projektlokale
Review Evidence behandelt und nicht als finaler Projektoutput ausgegeben.
Dadurch bleibt zwischen erzeugtem, validiertem, geprüftem und tatsächlich
veröffentlichtem Zustand eindeutig unterscheidbar.

Das veröffentlichte Output Package ist immutable und versioniert. Frühere
veröffentlichte Zustände werden bei späteren Änderungen nicht
überschrieben, sondern bleiben als eigenständige historische
Modellzustände nachvollziehbar.

Die Publication erzeugt selbst keine neue Engineering-Semantik. Sie
materialisiert ausschließlich den bereits validierten und durch den Human
Release autorisierten Modellzustand als finales Projektergebnis.

Damit endet der implementierte Engineering Information Flow des Proof of
Concept:

\begin{quote}
heterogene Engineering Sources\\
$\rightarrow$ source-grounded Processing\\
$\rightarrow$ Human Engineering Authority\\
$\rightarrow$ Project Fit\\
$\rightarrow$ Model Derivation und Human Placement\\
$\rightarrow$ Model Assembly Review und Base IEM\\
$\rightarrow$ Target-Model Formulation Authority\\
$\rightarrow$ Model-Quality Refinement und Human Model Quality Review\\
$\rightarrow$ human-authorized Successor IEM\\
$\rightarrow$ deterministische SysML-v2-Generierung\\
$\rightarrow$ technische Validierung\\
$\rightarrow$ Final Human Model Review\\
$\rightarrow$ kontrollierte Publication.
\end{quote}

Die nachfolgenden Abschnitte betrachten nicht mehr einzelne Schritte dieses
Verarbeitungswegs, sondern diejenigen technischen Realisierungsaspekte, die
mehrere Stufen der Verarbeitung gleichzeitig betreffen.

\subsection{Querschnittliche Realisierungsaspekte}
\label{sec:crosscutting_implementation}

Die vorangegangenen Abschnitte beschreiben die prototypische Realisierung
entlang des Engineering Information Flow. Mehrere technische Eigenschaften
lassen sich jedoch nicht ausschließlich einer einzelnen Verarbeitungsstufe
zuordnen.

Dies betrifft insbesondere die menschliche Interaktion mit dem System, die
durchgängige Bindung von Informationszuständen über Identitäten und
Fingerprints sowie die Beziehung zwischen der in Kapitel~4 beschriebenen
logischen Architektur und der konkreten Softwareimplementierung.

Diese Aspekte werden im Folgenden deshalb querschnittlich betrachtet.


\subsubsection{Human-facing Workflow und Decision-centered Interaction}
\label{sec:guided_workflow_impl}

Die umfangreiche interne Evidence-, Authority- und Traceability-Struktur des
Proof of Concept wird dem Nutzer nicht unmittelbar in Form der zugrunde
liegenden technischen Artefaktstruktur präsentiert.

Die Interaktion erfolgt über eine gemeinsame Streamlit-basierte Anwendung,
die den persistierten Engineering-Zustand in aufgabenorientierte
Arbeitsbereiche projiziert. Dazu gehören insbesondere Project- und
Source-Verwaltung, Processing, Human Engineering Review, Model Proposal beziehungsweise Model Placement,
Target-Model Formulation und Model Quality Review, Final Model Review und Published Output.

Die Benutzeroberfläche bildet dabei keine eigenständige Authority.
Engineering Decisions, Placement Decisions, Release Decisions und andere
autoritative Zustandsänderungen werden weiterhin durch die dafür vorgesehenen
Domain Services ausgeführt und persistent gespeichert. Die UI initiiert diese
Operationen lediglich auf explizit ausgewählten Artefaktzuständen.

Nach einer schreibenden Operation wird der sichtbare Zustand erneut aus den
persistierten Artefakten rekonstruiert. Session State oder Navigationszustand
der Oberfläche darf deshalb keine Engineering-, Model- oder Publication
Authority ersetzen.

Für die Darstellung wird das in ADR-017 formulierte Prinzip

\begin{quote}
\textit{Simple by default. Explainable on demand. Fully traceable underneath.}
\end{quote}

verwendet.

Im Standardworkflow steht daher zunächst die für die aktuelle Engineering-
Aufgabe relevante Information im Vordergrund. Technische IDs, Fingerprints,
Persistenzpfade und vollständige Provenance-Beziehungen bleiben verfügbar,
dominieren jedoch nicht die primäre Arbeitsoberfläche.

Die Oberfläche unterscheidet damit zwischen einer fokussierten Engineering-
Sicht und einer technisch detaillierteren Sicht. Beide Darstellungen basieren
auf demselben autoritativen Systemzustand. Der Wechsel der Darstellungstiefe
verändert keine persistierte Engineering Information.

Ein zweites Gestaltungsprinzip besteht in der
\emph{decision-centered interaction}. Offene menschliche Entscheidungen werden
als eigenständige Arbeitselemente sichtbar gemacht. Der Nutzer soll nicht erst
technische Processing-Zustände analysieren müssen, um festzustellen, an welcher
Stelle eine Entscheidung erforderlich ist.

Wo mehrere Persona-Perspektiven verfügbar sind, werden deshalb insbesondere
deren Vorschläge, Übereinstimmungen und relevante Abweichungen gemeinsam mit
der erforderlichen Human Decision dargestellt. Consensus und Variance dienen
dabei ausschließlich als Entscheidungsevidence. Auch visuell starke
Übereinstimmung darf nicht mit einer bereits erfolgten Human Approval
gleichgesetzt werden.

Das wiederkehrende Interaktionsmuster folgt damit grundsätzlich der Abfolge

\begin{quote}
Input\\
$\rightarrow$ Proposed Result\\
$\rightarrow$ Variance beziehungsweise relevante Unsicherheit\\
$\rightarrow$ Human Decision\\
$\rightarrow$ autorisierter Result State.
\end{quote}

Dieses Muster wird über mehrere Verarbeitungsebenen hinweg wiederverwendet.
Dadurch kann die interne technische Komplexität des Systems erhalten bleiben,
ohne für jede Authority Boundary ein vollständig anderes Interaktionsmodell
einzuführen.

ADR-024 konkretisiert diese Realisierung als Guided Engineering Workflow. Die
Guided-Workflow-Schicht bleibt dabei eine Projektion und Delegationsschicht über
den bereits vorhandenen Engineering Services und keine zusätzliche
Autoritätsinstanz.


\subsubsection{Provenance, Identitäten, Manifeste und Fingerprints}
\label{sec:provenance_identity_impl}

Ein zweiter querschnittlicher Realisierungsaspekt ist die durchgängige
technische Bindung der Informationszustände.

Die einzelnen Verarbeitungsschritte tauschen nicht ausschließlich fachliche
Inhalte aus. Relevante Artefakte tragen zusätzlich technische Identitäten,
Provenance-Informationen, Manifestdaten und Fingerprints, durch die ihr
konkreter Entstehungs- und Autorisierungszustand nachvollziehbar bleibt.

Die Identitäten der unterschiedlichen Artefakttypen werden bewusst nicht
vereinheitlicht. Eine Source, ein Processing Run, ein Engineering Subject,
eine Human Decision, ein Model Candidate, ein Internal Model Element und ein
Published Output repräsentieren unterschiedliche technische und fachliche
Konzepte und besitzen deshalb getrennte Identitäten.

Diese Trennung verhindert, dass semantische Kontinuität, technische
Artefaktidentität und Authority State implizit gleichgesetzt werden.

Manifeste bündeln die für einen Verarbeitungsschritt relevanten Referenzen und
Zustandsinformationen. Sie dokumentieren beispielsweise, welche Eingaben
konsumiert wurden, welche Artefakte daraus entstanden sind und welche
versionierten Profile oder Regelwerke für den jeweiligen Schritt maßgeblich
waren.

Fingerprints ergänzen diese logische Identifikation um die Bindung an einen
konkreten Inhalt. Eine Human Decision referenziert damit nicht nur ein
fachliches Subject, sondern den tatsächlich geprüften Informationszustand.
Entsprechend bindet ein Validation Result nicht lediglich einen Dateinamen,
sondern das exakte Generated SysML Artifact Set, auf dem die Validierung
durchgeführt wurde.

Dieses Prinzip setzt sich bis zur Publication fort. Ein Final Model Review
Decision kann nur den konkret geprüften Artifact- und Validation-Zustand
freigeben. Ändert sich einer dieser Zustände, darf die bestehende Entscheidung
nicht auf den neuen Inhalt übertragen werden.

Die Implementierung verwendet deshalb an Authority Boundaries überwiegend
immutable beziehungsweise append-orientierte Artefakte. Ein neuer Zustand
ersetzt einen früheren Zustand nicht stillschweigend, sondern wird als neuer
referenzierbarer Artefaktzustand persistiert.

Auch die Auswahl nachgelagerter Zustände erfolgt nicht über ein implizites
\enquote{latest wins}-Prinzip. Wo ein bestimmter Candidate Set, IEM Snapshot
oder Generated Artifact State konsumiert wird, muss dieser explizit
adressierbar und überprüfbar sein.

Dadurch entsteht über die gesamte Verarbeitung hinweg eine technische
Provenance-Kette von der ursprünglichen Source bis zum veröffentlichten
SysML-v2-Output. Diese Kette ist gleichzeitig die Grundlage dafür, menschliche
Authority Decisions an genau den Informationszustand zu binden, der tatsächlich
Gegenstand der jeweiligen Entscheidung war.


\subsubsection{Architektur-zu-Implementierung-Mapping}
\label{sec:architecture_implementation_mapping}

Die in Kapitel~4 entwickelte logische Architektur beschreibt
Verantwortlichkeiten und Informationsgrenzen des Turing Generators. Diese
Struktur wird in der Softwareimplementierung jedoch nicht als direkte
Eins-zu-eins-Abbildung der Logical Components auf Python-Module umgesetzt.

Mehrere technische Module können gemeinsam eine Architekturverantwortung
realisieren. Umgekehrt können technische Basisdienste, beispielsweise für
Persistenz, Identitätsbildung oder Fingerprint-Berechnung, von mehreren
Architekturverantwortlichkeiten genutzt werden.

Eine direkte Gleichsetzung von Logical Component und Softwaremodul würde daher
eine technische Struktur suggerieren, die weder für die Architektur noch für
die Implementierung maßgeblich ist.

Tabelle~\ref{tab:architecture_implementation_mapping} zeigt deshalb nur
repräsentative Zuordnungen zwischen den in Kapitel~4 beschriebenen
Architekturverantwortlichkeiten und ihrer prototypischen Umsetzung.

\begin{table}[htbp]
\centering
\caption{Repräsentatives Mapping zwischen Architekturverantwortlichkeiten und prototypischer Implementierung}
\label{tab:architecture_implementation_mapping}
\begin{tabular}{p{0.23\textwidth} p{0.27\textwidth} p{0.29\textwidth} p{0.13\textwidth}}
\hline
\textbf{Architekturverantwortung} &
\textbf{Repräsentative Implementierung} &
\textbf{Zentrale Artefakte und Contracts} &
\textbf{Design Rationale} \\
\hline

Project and Source Context Management &
Project-, Source- und Processing-Services &
Project, Source, Processing Run, Attempt, Manifeste &
ADR-005, ADR-012 \\

Engineering Information Processing &
Source Preparation, Evidence Detection, Engineering Subjects,
Persona Interpretation und Consensus &
Source Projection, Evidence, Engineering Subject,
Interpretation und Consensus Artifacts &
ADR-009, ADR-025 bis ADR-027 \\

Human Review and Engineering Authority &
Subject Review, Review Workspace und AEI Promotion &
Human Review Decision, Reviewed State,
Approved Engineering Information &
ADR-016 \\

Architecture Derivation and Model Placement &
Model Candidate Derivation und Model Placement Services &
Model Candidate Set, Placement Proposal,
Approved Model Placement Set &
ADR-018, ADR-029 \\

Model Assembly and Base Internal Model &
Model Assembly, Assembly Review und Internal Model Services &
Model Assembly Draft, Final Assembly Decision,
Base Internal Engineering Model &
ADR-019, ADR-029 \\

Target-Model Formulation and Model Quality Authority &
Target-Model Formulation, Model Refinement und Human Model Quality Review &
Formulation Proposals, Target-Model Formulation Authority,
Model Quality Authority, Successor IEM &
Target-Model-Formulation- und Model-Quality-Contracts \\

SysML v2 Artifact Generation &
Deterministische Generation Services und versionierte Generation Profiles &
Human-authorized Successor IEM,
Generated SysML Artifact Set und Traceability Entries &
ADR-021 \\

Generated Artifact Validation &
Interne Validatoren und SYSIDE Adapter &
SysMLValidationResult und Validation Findings &
ADR-022 \\

Final Model Review and Publication &
Final Model Review und Publication Services &
Final Model Review Decision und Published Output Package &
ADR-023 \\
\hline
\end{tabular}
\end{table}

Die Tabelle dient nicht als vollständige Softwaredokumentation. Insbesondere
werden technische Hilfsmodule und gemeinsam genutzte Infrastruktur bewusst
nicht vollständig aufgeführt.

Entscheidend ist vielmehr die Übereinstimmung der realisierten
Verantwortungsgrenzen mit der konzipierten Informationsverarbeitungsarchitektur.
Die logische Architektur beschreibt, welche Verantwortung getrennt werden
muss, während die konkrete Modulstruktur festlegt, wie diese Verantwortung im
Proof of Concept technisch umgesetzt wird.

Damit schließt sich die Verbindung zwischen der in Kapitel~4 entwickelten
Architektur und ihrer prototypischen Realisierung. Die anschließende
Verifikation und Validierung untersucht, inwieweit die implementierten
Mechanismen die vorgesehenen Informations-, Authority- und
Traceability-Grenzen tatsächlich einhalten.
